\section{Discussion}
\label{sec:discussion}

\todo{Write after results are finalized. Target: 2--3 pages.}

\subsection{Interpretation of Results}

\todo{Discuss the following points:
\begin{itemize}
  \item Why does PJM achieve the highest savings? Link to its high CV and coal/gas mix.
  \item Why does Singapore achieve the lowest savings? Low CV confirms the theoretical prediction.
  \item What does the temporal vs.\ spatial decomposition reveal? Which dimension dominates, and why does this vary across regions?
  \item Does the RF weighting ($\alpha > 0$) meaningfully change behavior vs.\ pure CI minimization?
\end{itemize}}

\subsection{Comparison with Prior Work}

\todo{Compare quantitative results with published benchmarks where possible:
\begin{itemize}
  \item \citet{wiesner2024limitations}: their empirical bounds on temporal shifting. Do our LP savings fall within their reported range?
  \item \citet{riepin2024spatiotemporal}: their CFE cost reduction per percentage of flexible load. Does our result corroborate their finding?
  \item \citet{xu2025green}: GPU cluster savings of up to 30\% from temporal shifting. How does our joint LP compare?
\end{itemize}}

The LP consistently outperforms the greedy baseline, demonstrating the value of joint optimization over rule-based heuristics.
\todo{Quantify the gap once results are available.}

\subsection{Limitations}

\begin{enumerate}[leftmargin=1.5em]
  \item \textbf{Perfect historical information.}
        The LP is evaluated in backtesting mode with full knowledge of the 2024--2025 CI/RF series.
        Real-time deployment requires CI forecasting; forecast error will reduce achievable savings, particularly in high-variability grids like PJM.

  \item \textbf{Homogeneous workload model.}
        All flexible load is treated as identical.
        Heterogeneous job types (training vs.\ inference vs.\ preprocessing) with different delay tolerances, power profiles, and checkpointing costs are not modeled.

  \item \textbf{Stationary demand.}
        Total demand $D$ is fixed; demand elasticity, dynamic arrivals, and load shedding are not considered.

  \item \textbf{PUE abstraction.}
        Regional differences in Power Usage Effectiveness are abstracted away.
        Incorporating PUE would amplify or attenuate the carbon savings estimates for each region.

  \item \textbf{ElectricityMaps data provenance.}
        CI and RF values are ElectricityMaps' modeled outputs, not raw grid measurements.
        Modeling error in the data source propagates into the optimization solution but cannot be separately quantified without access to raw TSO data.
\end{enumerate}
